\clearpage\thispagestyle{empty}
\pdfbookmark[0]{Quellen zum Pallas-Bericht vom April 1805}{pallas180504-sources}
\begin{center}\textbf{Quellen zum Pallas-Bericht vom April 1805}\end{center}
\noindent Diese KI-Übertragung in heutiges Deutsch folgt dem Auszug aus Franz Xaver von Zachs Bericht in \href{https://gdz.sub.uni-goettingen.de/id/PPN236007467}{Gauß’ \textit{Werken}, Band~VI} (Göttingen, 1874), S.~262--263. Die entsprechenden Passagen aus Beitrag XXXIII, \textit{Fortgesetzte Nachrichten über den neuen Haupt-Planeten Pallas}, erschienen im Aprilheft der \textit{Monatlichen Correspondenz}, Band~XI, auf \href{https://www.digitale-sammlungen.de/en/view/bsb10538603?page=388}{S.~376}, \href{https://www.digitale-sammlungen.de/en/view/bsb10538603?page=389}{S.~377}, \href{https://www.digitale-sammlungen.de/en/view/bsb10538603?page=393}{S.~381} und \href{https://www.digitale-sammlungen.de/en/view/bsb10538603?page=395}{S.~383}.
\par\smallskip\noindent Von Zach berichtet über die achten Elemente von Carl Friedrich Gauß, nennt die Beobachter und beschreibt Karl Ludwig Hardings Arbeit an einer Sternkarte. Die Ephemeride steht in der Zeitschrift auf \href{https://www.digitale-sammlungen.de/en/view/bsb10538603?page=389}{S.~377} bis \href{https://www.digitale-sammlungen.de/en/view/bsb10538603?page=392}{S.~380}. Die Werkausgabe beschreibt sie in eckigen Klammern und verweist auf den \hyperlink{native-267}{Ceres-Bericht auf S.~262}.
\par\smallskip\noindent Übertragung und Satz wurden mit KI erstellt. Die Digitalisate der Zeitschrift stammen von der Bayerischen Staatsbibliothek.
\clearpage
